\chapter{Netzwerk}

\section{Motivation und Testumgebung}
Die Implementierung eines Mehrspieler-Szenarios stellte eine zentrale Herausforderung bei der Entwicklung des Spiels dar.
Zwar verfügt \texttt{MOROS} über einen rudimentären Netzwerkstack, die in \texttt{QEMU} integrierte \texttt{listen/connect}-Funktionalität erlaubt jedoch nur Punkt-zu-Punkt-Verbindungen.
Damit war ein echtes Multiplayer-Setup, bei dem mehrere Clients gleichzeitig mit einem Server interagieren, nicht realisierbar.
Um dieses Problem zu umgehen, kam das \emph{Virtual Distributed Ethernet (VDE)} zum Einsatz.
VDE ermöglicht die Einrichtung virtueller Switches und damit die Verbindung mehrerer virtueller Maschinen über ein gemeinsames virtuelles Netzwerk, was insbesondere für Mehrspieler-Szenarien von essenzieller Bedeutung ist.

\section{Server-Client-Architektur}
Die gewählte Architektur folgt dem klassischen Server-Client-Paradigma.
Der Server ist für die zentrale Berechnung des Spielzustands verantwortlich, während die Clients lediglich Eingaben liefern und den aktuellen Zustand rendern.

\subsection{Ablauf im Server}
Der Server sammelt zunächst alle Eingaben der verbundenen Clients ein.
Anschließend wird der globale Spielzustand berechnet, indem Bewegungen der Spieler, Kugeltrajektorien und Kollisionen aktualisiert werden.
Am Ende eines jeden Ticks sendet der Server den konsistenten Spielzustand in Form eines Broadcasts an alle Clients zurück.

\subsection{Ablauf im Client}
Der Client ist primär für zwei Aufgaben verantwortlich: zum einen für das Senden eigener Eingaben (Tastatur- und Mausevents), zum anderen für das Rendern des vom Server empfangenen Spielzustands.
Die Spiel-Logik selbst wird weitgehend serverseitig berechnet, um Inkonsistenzen zu vermeiden.
Dennoch existieren bestimmte Mechanismen, bei denen der Client selbst Berechnungen übernimmt, um die Netzwerklast zu reduzieren (z.\,B. bei der Pfadberechnung von Projektilen).

\section{Nachrichtenformate und Protokolle}
Die Kommunikation zwischen Server und Clients basiert auf einem kompakten binären Protokoll.
Zentrale Nachrichtentypen sind:

\begin{itemize}
    \item \texttt{Connect}: Aufbau einer neuen Verbindung
    \item \texttt{MapRequest}/\texttt{MapData}: Anforderung und Übertragung der aktuellen Karte
    \item \texttt{PlayerInput}: Übermittlung der Steuerungsinformationen des Clients
    \item \texttt{GameStateUpdate}: Aktualisierter Spielzustand vom Server
    \item \texttt{Disconnect}: Abmeldung eines Clients
\end{itemize}

Um die Übertragung effizient zu gestalten, werden Eingaben in Bitmasken codiert.
Richtungsangaben und Schussaktionen lassen sich auf ein einzelnes Byte reduzieren, wodurch der Overhead pro Client signifikant verringert wird.
Komplexere Objekte wie Kugeln werden nicht vollständig übertragen; stattdessen erhalten Clients nur Startpunkt, Richtung und Spieler-ID.
Die exakte Pfadberechnung erfolgt lokal, was die Paketgröße drastisch reduziert (von über 100 Bytes auf lediglich 24 Bytes pro Projektil).

\section{Synchronisation und Robustheit}
Ein zentrales Problem in verteilten Spielen ist der Umgang mit asynchronen Zuständen und Paketverlusten.
Die gewählte Lösung basiert auf periodischen State-Updates des Servers, die sicherstellen, dass Clients regelmäßig synchronisiert werden.
Kurzfristige Paketverluste können so kompensiert werden, da der nächste gültige Spielzustand überschreibend angewandt wird.

Als Transportprotokoll kommt \texttt{UDP} zum Einsatz.
Dieses bietet zwar keine garantierte Zustellung, ist jedoch aufgrund seiner geringen Latenz und des reduzierten Overheads für Echtzeitanwendungen prädestiniert.
Einschränkungen ergeben sich aus der maximal sicheren \emph{Maximum Transmission Unit (MTU)} von etwa 1500 Bytes.
Werden Pakete größer, erfolgt eine Fragmentierung auf IP-Ebene, die bei Verlust einzelner Fragmente zum Ausfall des gesamten Pakets führt.
In der Praxis limitiert dies die maximal mögliche Anzahl an Spielern und gleichzeitig existierenden Projektilen pro Spielrunde.

